% Generated draft: requires mathematical and source review.
\documentclass[11pt]{article}
\usepackage[margin=25mm]{geometry}
\usepackage{fontspec}
\setmainfont{Latin Modern Roman}
\usepackage{amsmath,amssymb,mathtools,unicode-math}
\setmathfont{Latin Modern Math}
\usepackage{enumitem,needspace,xurl,hyperref}
\setlength{\parindent}{0pt}
\setlength{\parskip}{5pt}
\setlength{\emergencystretch}{3em}
\newcommand{\N}{\mathbb{N}}
\newcommand{\Z}{\mathbb{Z}}
\newcommand{\Q}{\mathbb{Q}}
\newcommand{\R}{\mathbb{R}}
\newcommand{\C}{\mathbb{C}}
\newcommand{\F}{\mathbb{F}}
\newcommand{\A}{\mathbb{A}}
\newcommand{\PP}{\mathbb{P}}
\newcommand{\E}{\mathbb{E}}
\newcommand{\eps}{\varepsilon}
\newcommand{\dd}{\,\mathrm{d}}
\begin{document}
\section*{UnsolvedMath ID: 30000547}
\textbf{Source problem code:} catalogue entry
\subsection{Definitions and mathematical statement}
Let $\Gamma<\operatorname{PU}(2,1)$ be the uniform arithmetic lattice in the source, generated by complex reflections of order five. Uniform means that $\Gamma\backslash\mathbb H^2_{\C}$ is compact. A complex reflection fixes a complex geodesic pointwise. Let $\Sigma_g$ denote a closed orientable surface of genus $g$ and put $H=\pi_1(\Sigma_g)$.

The requested map is written literally as
\[
F:\mathbb H^2_{\R}\times H\longrightarrow\mathbb H^2_{\C}.
\]
With $H$ discrete, the domain is a disjoint union of real hyperbolic planes. Under this interpretation, ``totally geodesic real embedding'' means that $F$ is an embedding and each restriction $F_h(x)=F(x,h)$ embeds a totally geodesic real hyperbolic plane, with a metric normalization compatible with the complex hyperbolic metric. In particular, different components must have disjoint images for the whole map to be injective.

For the intended action $\alpha$ of $\Gamma$ on this domain and the intended set $\mathcal R$ of reflections, the existence question requires
\[
F(\alpha(\gamma,z))=\gamma F(z)
\quad(\gamma\in\Gamma),\qquad
F(\mathbb H^2_{\R}\times H)\cap
\bigcup_{r\in\mathcal R}\operatorname{Fix}(r)=\varnothing.
\]
The extract does not define $\alpha$, $\mathcal R$, or $g$. Nor does it identify a particular lattice by parameters or generators, although the literature assessment calls it a Deligne--Mostow lattice. Thus the product-domain interpretation and displayed conditions explain the written target but do not supply the missing equivariant structure. The factor $\pi_1(\Sigma_g)$ is not silently replaced by a quotient, a surface, or a homogeneous coset set.
\subsection{Short English statement}
Ask for an equivariant family of embedded totally geodesic real hyperbolic planes in the complex hyperbolic plane, disjoint from the relevant reflection mirrors. The domain action, genus, precise lattice and reflection collection are not supplied, so the literal product notation remains qualified.
\subsection{Sources}
\begin{itemize}
\item Supplied report: \emph{Geometric Group Theory, Hyperbolic Dynamics and Symplectic Geometry} (2007), Oberwolfach Reports, DOI 10.4171/owr/2006/33, pp. 1991--2058; proposed year 2006. \url{https://doi.org/10.4171/owr/2006/33}.
\item Local UnsolvedMath registry, permanent ID 30000547, catalogue entry; created 25 August 2026, updated 16 September 2026. The corrected\_verified note describes meaning-preserving cleanup. Its supplied 21 August 2026 literature triage labels the specific Deligne--Mostow lattice problem open and says available reflection or surface-subgroup results do not resolve this equivariant embedding question. No additional external reference is supplied. The dated assessment has not been independently checked here.
\end{itemize}
\end{document}
